\begin{afterword}
\summaryhead

The post describes how people come to press advice on others. You read nineteen pages of self-improvement advice, try two methods that fail, and then find a twentieth that works. Now you know ``the real way,'' and you tell a friend.

The author is often on the receiving end. A board member of the author's institute once said the author needed no raise, because an online coupon service would cut food costs. The author believed it because a trusted friend said it; the coupons were given up after two weeks. The lesson: personal advice is no more likely to work than a blog post by any intelligent writer. The range of advice the author receives suggests that people differ in what works for them.

``Different things work for different people'' can be an evasion, the author grants, but until the general laws are known, the adviser should accept ``No.'' This matters most for anyone with power or leverage, since power is easy to abuse unnoticed, people underestimate how different others are, and few can tell laziness from real limits. The author resolves to accept ``No'' and to keep looking for deeper laws.

\responsehead

The post is a caution drawn from the author's experience, and it claims little more than that. Its central sentence, that a friend's earnest advice ``is no more and no less likely to work for me than a random personal improvement blog post,'' is about advice given to the author, and its support is the author's history as a frequent target of such advice. The exact equality is not measured. It works as a rule of thumb, and the post treats it as one: it grants that ``Sometimes it works,'' that ``Sometimes'' people are lazy, and that the author does not know ``when and why and for who'' a method works.

The post answers the strongest objection itself. ``Different things work for different people'' is a phrase often used to shield a claim from criticism, and the post says so. Its reply is that the phrase is a statement of ignorance, to be replaced by general laws when they are found. The next post in the book, ``Practical Advice Backed By Deep Theories,'' takes up that search.

Its general claim, ``We underestimate the distance between ourselves and others,'' comes with no source, but research supports it. The false consensus effect, named by Ross, Greene and House in 1977, is the finding that people see their own choices as common and appropriate, and make more extreme judgments about those who choose differently. That fits the adviser who reads another person's failure as laziness.

The opening anecdote already connects the two halves of the post, since the coupon tip came from a board member and was offered as a reason against a raise. In a comment the same day, the author explained why the post belongs in this sequence on communities: other-optimizing is ``a mode by which managers sabotage the managed.''

In short: a caution drawn from the author's experience and offered as a rule of thumb; it answers its own strongest objection, and research on false consensus supports its claim that we underestimate how different others are.
\end{afterword}
